\chapter{没有电路图，怎么修机器}
% 完整版：chapters/17_debugging.tex

\pencilfig{17}{\pencilart{17}}{没有电路图的电路板：探针在数十亿根导线中只听得到几根。}

想象有人给你一块正在工作、却没有电路图的电路板，请你把它修好。工程师有四样工具：探针，用来听信号；信号发生器，用来输入自己的信号；拆掉一个零件、看看什么坏了的能力；以及访问控制寄存器的途径。研究大脑用的正是这几样工具。

\section*{听一千个神经元的探针}

插进大脑的电极能听到附近几个神经元的咔嗒。看起来，相对于八百六十亿个神经元，一千个电极微不足道。但它们足以控制屏幕上的光标。为什么？相邻神经元的噪声很相似，而控制手臂的活动实际上只需要一二十个共同变量就能描述。不必听所有神经元，只要听到这几个“声音”就够了。

一千个电极的原始信号每秒有几亿比特：从颅骨里发出的无线信道传不了这么多。而咔嗒本身携带的信息大约只有它的千分之一。所以，在植入物上直接提取出咔嗒、只把它们发送出去是划算的——和眼睛里用的是同一种压缩手法。

\section*{Neuralink 在做什么}

Neuralink 公司着手解决这个问题的工程部分：用细而柔软的丝线代替硬针，用机器人把它们植入，再加上密封的无线植入物。在它已发表的论文中，描述的是一套用电缆连接计算机的系统，压缩和无线通信被列为计划。关于一位双手瘫痪的人控制光标的数据，目前主要来自公司自己的报告；据这些报告，植入物大约有一千个通道，分布在几十根丝线上，部分丝线在植入后发生了移位。这个想法本身并不新：学术团队使用约有一百个电极的硬质阵列，早已让人控制光标、用意念“写”字母，并把说话的尝试以大约每分钟六十个词的速度变成屏幕上的文字。

\section*{读取与写入}

解码器读取许多神经元的咔嗒频率，得到的信息每秒不到十比特——只是神经元本身所携带信息中微不足道的一小部分。好在大脑和解码器会互相学习：人适应解码器，解码器也适应人。两个彼此适应的学生是一个微妙的问题：即使每一个单独都是稳定的，合在一起也可能摇摆起来。

向大脑写入比读取更难。电极上的电流会激发周围的一切，而不是只激发需要的那个神经元。可以“画”出光点，就像发光的像素，人能分辨出简单的字母。但要写入眼睛的压缩编码、让大脑看见一幅图像，目前还没有人做得到。

\section*{探针看不见什么}

电极听到的是数据电话网。而电台——多巴胺、血清素、去甲肾上腺素、乙酰胆碱——是物质的浓度，探针听不到它们。它也看不见存储着全部记忆的连接强度。它同样看不见人所感受到的那种疼痛。一个看得见数据、却看不见寄存器的调试器，永远会困惑：为什么同样的输入会给出不同的输出。

\fullversion{17}
